\section{Related Work}

\paragraph{Execution Traces for Code LLMs.}
TracePile \citep{chen2025tracepile} scaled trace-augmented training to 2.6M samples; \citet{wang2025dcsh} found trace utility task- and format-dependent, which we extend by identifying exception events as the dominant signal.
NExT \citep{ni2024next} and ET-CoT \citep{eguchi2025etcot} used trace chain-of-thought for repair and execution reasoning; \citet{ferrando2025execute} found compact traces comparable to full traces, supporting our density finding.
\citet{liu2023codeexecution} pre-trained on execution traces; CRUXEval \citep{gu2024cruxeval}, extended to 19 languages \citep{xu2025cruxevalx}, established execution reasoning benchmarks.
None decompose which trace components contribute to \emph{classification} accuracy.

\paragraph{Code Verdict and Patch Classification.}
TRACED \citep{ding2024traced} showed execution-aware pretraining improves vulnerability detection; DeepDFA \citep{steenhoek2024deepdfa} and learned program slicing \citep{yadavally2024slicing} combined classical analysis with deep learning.
In patch correctness, \citet{tian2020patchcorrectness}, CC2Vec \citep{hoang2020cc2vec}, LLM4PatchCorrect \citep{zhou2024llm4patch} (84.4\%), and APPT \citep{appt2024} (79.7\%) developed representation learning of code changes.
Verdict prediction is an instance of the longstanding test oracle problem \citep{barr2015oracle}; FuseFL \citep{widyasari2024fusefl} and \citet{foster2025mutation} demonstrated structured fault signals; the equivalent-mutant problem \citep{papadakis2019mutation} poses an analogous classification challenge.
None isolate individual trace components or compare against rule-based baselines.

\paragraph{Fault Localization.}
Spectrum-based fault localization (SBFL) ranks program elements by suspiciousness from test spectra \citep{abreu2007accuracy,jones2002visualization,wong2014dstar}; recent LLM-based approaches \citep{kang2024autofl,qin2024agentfl} share our insight that anomalous execution events carry the strongest diagnostic signal, with verdict prediction serving as a prerequisite to FL.

\paragraph{Trace Compression and Selection.}
TraceCoder \citep{huang2026tracecoder} and InspectCoder \citep{wang2025inspectcoder} use selective trace extraction; RepairAgent \citep{bouzenia2025repairagent}, Agentless \citep{xia2024agentless}, SWE-agent \citep{yang2024sweagent}, and OpenHands \citep{wang2025openhands} incorporate execution feedback in repair workflows \citep{jimenez2024swebench}, yet which trace components carry verdict signal most efficiently remains unstudied.
LLMLingua \citep{jiang2023llmlingua} and Selective Context \citep{li2023selectivecontext} showed that pruning low-information tokens preserves accuracy; \citet{haque2025traces} found raw traces provide limited repair benefit.
Our decomposition (${\sim}$24\% gradient dilution vs.\ ${\sim}$76\% information overload) provides the first quantitative framework for \emph{why} compression helps.

\paragraph{Calibration and Defect Prediction.}
Neural network calibration \citep{guo2017calibration,kadavath2022language,spiess2025calibration} studies confidence-correctness alignment; our ECE ${<}$ 0.007 suggests compact inputs improve calibration. Just-in-time defect prediction \citep{kamei2013large,hoang2019deepjit} shares our fault detection motivation but operates on code change features.

\paragraph{Training Efficiency and Data Selection.}
Curriculum learning \citep{bengio2009curriculum} and data selection methods \citep{swayamdipta2020cartography,sachdeva2024dataefficient} show that data quality improves convergence, consistent with our finding that compact exceptions converge rapidly while verbose traces oscillate.
